\section{Preregistered Audit Design}
\label{sec:design}

\paragraph{Models, data and roles}
The audit, registered in a development stage and a validation stage whose
results were then frozen, uses OpenAI CLIP ViT-B/32 and ViT-L/14 with frozen weights and
frozen prompts: the single OpenAI CIFAR-100 prompt and a three-template EuroSAT
ensemble~\cite{radford2021clip,krizhevsky2009cifar,helber2019eurosat}. From the
CIFAR-100 training split and the EuroSAT image folder we drew three folds, each
with five development and five validation images per class, giving
\ExpCifarItemsPerCell{} validation items per CIFAR-100 cell and
\ExpEurosatItemsPerCell{} per EuroSAT cell, \ExpCells{} cells and
\ExpPositions{} validation positions in total. Development items served only to
build operators, and validation items to select and evaluate them. The
CIFAR-100 test split and a EuroSAT holdout were sealed and were not read by the
audit, the positive control or the extension, whose CIFAR-10 items also come
from the training split; the
registered follow-up of Section~\ref{sec:followup} later spent the sealed
EuroSAT holdout (\FUItemsEurosatB{} items) and \FUItemsCifarB{} CIFAR-100
official-test items once, as its registration specifies.

\paragraph{Certification}
Every candidate bank in every cell was certified on every validation item with
independent selection and confirmation noise streams derived from separate base
seeds, so all \ExpCandidates{} banks of a cell see identical draws (a paired
design), while selection draws are never reused for confirmation. In total the
audit evaluated \ExpRows{} candidate--item positions with \ExpSelVotes{}
selection votes and \ExpConfVotes{} confirmation votes. Radii are reported at
$r\in\{0,0.1,0.25,0.5\}$ with $r=0.25$ preregistered as primary.

\paragraph{Selection and efficacy rule}
The protocol selected one proposal objective and step by the macro
\emph{anchored} certified accuracy at $r=0.25$ over the \ExpCells{} cells, with
frozen tie-breaks, and selected one coefficient per tunable baseline by the same
metric. A candidate was eligible only if its macro clean-accuracy drop relative
to the uncorrected model was at most one point and no single cell dropped by
more than two points. The efficacy gate then required the selected proposal to
exceed the eligible mean-centering comparator on the macro metric and on every
model and dataset marginal. If the gate failed, the protocol stopped: the
preregistered attack and control stage and the sealed test remained
locked, and the project committed to reporting the negative in full. That
stage never ran; the sealed test was later allocated, by a registration hashed
before sampling, to the follow-up
(Section~\ref{sec:followup}).

\paragraph{Custody}
Configurations, split assignments, seeds, prompts and every bank were hashed
before sampling; the reconstruction checks are in
Appendix~\ref{app:deviations}.

\paragraph{Post-hoc audits of the immutable result}
Three analyses were specified after the outcome was known and are labeled as
such. (i) \emph{Metric semantics}: standard certified accuracy
\eqref{eq:stdca} was reconstructed beside the anchored metric \eqref{eq:anchca}
for all candidates, cells and radii. (ii) \emph{Complete-family bands}: a paired
item bootstrap, stratified by dataset, fold and class, resampled the same items
across both models, all candidates, radii and metrics, and produced one
unstudentized 95\% simultaneous max-absolute-deviation band over all
\BandContrasts{} candidate--radius--metric contrasts against the uncorrected
model (\BandReplicates{} replicates, \BandStrata{} strata, seed fixed before
execution). The band is a conditional compatibility statement about the saved
sample; it is not an equivalence test. (iii) \emph{Exact-count adjudication}:
the clean gate was recomputed with integer counts because the implementation
compared floating-point percentages.
